
\begin{abstract}
We introduce \emph{refinement reflection}, a method
to extend \emph{legacy} languages---with highly tuned
libraries, compilers, and run-times---into theorem provers.
%
The key idea is to reflect the code implementing a
user-defined function into the function's (output)
refinement type.
%
As a consequence, at \emph{uses} of the function,
the function definition is unfolded into the refinement logic
in a precise and predictable manner.
%
We have implemented our approach in \toolname
thereby retrofitting theorem proving into Haskell.
%
We show how to use reflection to verify
that many widely used instances of the Monoid,
Applicative, Functor, and Monad typeclasses
actually satisfy key algebraic laws required to 
make the code using the typeclasses safe.
%
Finally, transforming a mature language---with
highly tuned parallel runtime---into a theorem
prover  enables us to build the first
\emph{deterministic parallelism library}
that verifies assumptions about associativity
and ordering---that are crucial for determinism
but simply assumed by existing systems.
\end{abstract}

% \vspace{-5mm}

%%\begin{abstract}
%%Refinement Reflection turns your favorite programming
%%language into a proof assistant by reflecting
%%the ​code implementing a​ user-defined function
%%into the function's (output) refinement type.
%%%
%%As a consequence, at \emph{uses} of the function,
%%the function definition is unfolded into the refinement logic
%%in a precise, predictable and most
%%importantly, programmer controllable way.
%%%
%%In the logic, we encode functions and lambdas using uninterpreted symbols
%%preserving SMT-based decidable verification.
%%In the language, we provide a library of combinators that lets programmers
%%compose proofs from basic refinements and function definitions.
%%%
%%We have implemented our approach in the Liquid Haskell system,
%%thereby converting Haskell into an interactive proof assistant,
%%that we used to verify a variety of properties ranging
%%from arithmetic properties of higher order, recursive functions
%%to the Monoid, Applicative, Functor and Monad type class laws
%%for a variety of instances.
%%%
%%\new{Finally, transforming a mature language---with a parallel runtime---into a proof
  %%assistant enables us to build the first {\em deterministic parallelism
    %%library} that verifies assumptions about associativity and
  %%ordering---necessary to determinism but simply assumed by existing systems.}
%%\end{abstract}
